%%% introduction

\section{Introduction}
\label{sec:intro}

Solar system small bodies are relatively small objects orbiting around the Sun, such as asteroids, comets, and meteoroids. More than one million objects have been discovered as of August 2020\footnote{The number is adopted from the Minor Planet Center (\today).}. The objects whose perihelion distances are smaller than 1.3 au are called Near-Earth Objects (NEOs). Their population, orbits, sizes, and compositions provide vital information to understanding the solar system's evolution \citep[e.g.,][]{granvik_debiased_2018,bottke_debiased_2002}. Besides, some of them are potentially hazardous to the Earth. Thus, detecting NEOs is crucial as well in the context of planetary defense. However, many NEOs, especially smaller ones, remain undiscovered \citep[e.g.,][]{harris_population_2015}.

Multiple images of the same field are obtained with some intervals to detect moving objects in optical observations. Moving objects should appear in different positions of the images. A short arc is obtained by connecting such two sources, usually referred to as a tracklet. Connecting multiple tracklets provides a longer track, leading to a more precise orbit determination of a moving object.

Algorithms connecting tracklets have been developed so far. Moving Object Processing System \citep[MOPS;][]{denneau_pan-starrs_2013,kubica_efficient_2007} is the algorithm used in Pan-STARRS, utilizing a tree-based structure to efficiently connect inter-night tracklets. ZTF's Moving Object Discovery Engine \citep[ZMODE;][]{masci_zwicky_2019} is another implementation for Zwicky Transient Facility. ZMODE constructs a "stringlet" \citep{waszczak_main-belt_2013} from three or more detections, making the system more robust. HelioLinc provides a new idea to this field, where tracklets are connected in a heliocentric frame instead of a geocentric one. Although some algorithms are independent of tracklets \citep[THOR;][]{moeyens_thor_2021}, extracting tracklets is the initial step in moving object discovery.

The cadence of survey observations has increased with advances in observational instruments \citep[e.g., Pan-STARRS, ZFS, ATLAS, ASAS-SN, and EvryScope;][]{chambers_pan-starrs1_2016,magnier_pan-starrs_2016,dekany_zwicky_2020,tonry_atlas:_2018,kochanek_all-sky_2017,law_evryscope_2015}. Video observation is an extreme case \citep[e.g., Tomo-e Gozen, OASES, and TAOS2;][]{sako_tomo-e_2018,arimatsu_organized_2017,huang_taos_2021,lehner_status_2018}. An efficient data reduction process is required due to the enormous amount of data in the video observation. Parallel computing is one solution. \citet{yanagisawa_automatic_2005} developed an FPGA board to extract faint tracklets from a sequence of images. This paper presents another solution, developing an efficient algorithm. The proposed algorithm extracts linearly aligned points in a three-dimensional space, a fundamental routine to extract tracklets from video data.

The paper is organized as follows. The design and requirements of the proposed algorithm are described in Section 2. Then, the performances of the algorithm are illustrated in Section 3. Section 4 summarizes the paper. The algorithm we develop is available as open-source software.

%%% Local Variables:
%%% mode: latex
%%% TeX-master: "../manuscript"
%%% End:
